\section{An exact simple eigenfunction after a small perturbation}
\label{sec:exact}

Fix the input pair, \(l_-,l_+\), and put
\(w_{\rm c}=l_+-l_->0\),
\(s_i=l_-+i w_{\rm c}/5\) for \(1\le i\le4\).
For the allowed \(k=\sqrt{N(N+3)}\), let
\[
 b_N(t)=\frac{N}{2k}\log(1-t).
\]
Choose the additive constant in the envelope so that
\(h(s_2)=b_N(s_2)\). It is bounded uniformly in \(N\).
For the fixed small \(a_*\), and all sufficiently large \(N\),
\[
 \partial_t(h-b_N)
       =-a_*+\frac{N}{2k(1-t)}\ge c_0>0
       \quad(0\le t\le1/2).
\]
It follows that
\begin{equation}\label{eq:gluegaps}
 h\le b_N-c w_{\rm c}\ (t\le s_1),\qquad
 h\ge b_N+c w_{\rm c}\ (s_3\le t\le s_4).
\end{equation}
Take \(0\le\beta\le1\), equal to one for \(t\le s_3\) and
zero for \(t\ge s_4\). The smooth positive weight
\begin{equation}\label{eq:weight}
 W=e^{kh}+\beta(t)e^{kb_N(t)}
\end{equation}
uses a zero extension of the second summand.
The pair is round on the support of \(\beta\).
Thus \((L+k^2)(\beta u_*)\) is only a cutoff commutator.
Its support lies in the exponential gap of \eqref{eq:gluegaps}.
For every fixed \(q\ge0\) and \(M>0\),
\begin{equation}\label{eq:beamweight}
 |D^q(\beta u_*)|\le C_qk^qW,\quad
 |D^q(L+k^2)(\beta u_*)|\le C_{q,M}k^{-M}W,\quad
 |\dd W|/W\le Ck^2 .
\end{equation}
Derivatives of the fixed collar cutoff can depend on \(w_{\rm c}\);
the positive gap absorbs them when \(k\) is sufficiently large.

\subsection{A real sample with a nonvanishing first jet}

Choose a net in \(\{t\ge s_1\}\) of mesh at most \(k^{-10}\).
There are at most \(Ck^{40}\) centers. At each take the four
packets of Proposition~\ref{prop:packets}, and denote the collection
by \(\mathcal W_k\). Use independent standard real Gaussian
coefficients \(g_v^{(1)},g_v^{(2)}\) and set
\begin{equation}\label{eq:random}
 u=\beta u_*+\sum_{v\in\mathcal W_k}
       (g_v^{(1)}\operatorname{Re}v+
                         g_v^{(2)}\operatorname{Im}v).
\end{equation}
Let \(\mathcal B_k\) be the event that all coefficients have
absolute value at most \(k\). The elementary Gaussian tail bound gives
\(\Prob(\mathcal B_k^c)\le Ck^{40}e^{-k^2/2}\to0\).

Given finite \(J\ge2\) and \(M>0\), choose the packet residual
orders through \(J\) to exceed \(M+42\). On \(\mathcal B_k\),
their number, \eqref{eq:packetder} and \eqref{eq:beamweight} give
\begin{equation}\label{eq:samplebounds}
 |D^q u|\le C_qk^{2q+41}W,\qquad
 |D^q\mathfrak r|\le C_qk^{-M}W\quad(0\le q\le J),\qquad
 \mathfrak r=\Div_\mu(A\,\dd u)+k^2\rho u.
\end{equation}
All orders are chosen before \(k\).
Every packet vanishes for
\(t<s_1-Ck^{-1/3-100\delta}\), so in that region \(u=u_*\)
and \(\mathfrak r=0\).

\begin{lemma}\label{lem:jet}
For every fixed choice of packet orders,
\[
 \Prob\left(\mathcal B_k\cap
   \left\{\inf_B\frac{|u|+|\dd u|_{g_0}}{W}\ge k^{-500}\right\}\right)
 \longrightarrow1 .
\]
\end{lemma}

\begin{proof}
Where \(W/e^{kh}>k^{50}\), the cutoff is one for large \(k\),
by \eqref{eq:gluegaps}. The cosine and angular derivative of
the uncut beam give
\(|u_*|+|\dd u_*|_{g_0}\ge ce^{kb_N}\).
On \(\mathcal B_k\) the packet first jets are at most
\(Ck^{43}e^{kh}\), so the beam dominates there and gives a
lower bound \(cW\). This region includes \(t<s_1\).

In the remaining region there is a nearby center. The real
jet spanning from Proposition~\ref{prop:packets}, converted
to ordinary derivatives, gives a covariance matrix for the
five-dimensional value--gradient vector whose least eigenvalue
has square root at least
\[
 ck^{-1}e^{kh}\ge ck^{-51}W\ge k^{-52}W.
\]
One packet family suffices; others add nonnegative covariance.
The deterministic beam shifts only the mean.
Hence the density of the normalized vector
\(\mathcal J=W^{-1}(u,\partial_1u,\ldots,\partial_4u)\)
is at most \(Ck^{260}\). For fixed \(C_0,D>0\),
\begin{equation}\label{eq:smallball}
 \Prob(|\mathcal J(x)|\le C_0k^{-D})\le Ck^{260-5D}.
\end{equation}
This is an unconditional Gaussian estimate.

On \(\mathcal B_k\), two derivatives in
\eqref{eq:samplebounds} and \(|\dd W|/W\le Ck^2\) give
\(|D\mathcal J|\le Ck^{45}\le k^{46}\).
In a finite atlas take grids of mesh \(k^{-D-48}\), with
at most \(Ck^{4D+192}\) points. Apply \eqref{eq:smallball}
at the Gaussian-region grid points and the beam bound elsewhere.
The union bound is \(Ck^{452-D}\). With \(D=500\) it tends
to zero, while the interpolation error is \(O(k^{-D-2})\).
Choose \(C_0\) to cover the fixed coordinate norm comparisons.
This yields the stated global first-jet bound, after also
subtracting \(\Prob(\mathcal B_k^c)\).
The use of \(\mathcal B_k\) is confined to derivative control
and interpolation, not to the Gaussian density calculation.
\end{proof}

\subsection{Many strict signs at the local wavelength}

\begin{lemma}\label{lem:signs}
For fixed packet orders there are \(p_*,c_*>0\) such that,
for all sufficiently large allowed \(k\), \(u\) has a sign
certificate of size at least \(c_*k^{1+\gamma}\) with
probability at least \(p_*\). Potential disks can be fixed
deterministically before sampling.
\end{lemma}

\begin{proof}
In \(V\), use radii \(R_x=(kQ_x)^{-1}\).
Since \(|D Q|\le Ck^{4\delta}\), intersecting balls have
comparable \(Q\)'s and radii for large \(k\).
At a fixed \(k\) these radii have a positive minimum.
Choose a finite maximal disjoint family
\(\{B(x,R_x):x\in\mathcal X_k\}\).
Fixed enlargements cover \(V\), and \(Q\) is comparable
to \(Q_x\) in each enlargement. Therefore
\begin{equation}\label{eq:packing}
 \int_V kQ\,\dd x\le C\sum_{x\in\mathcal X_k}
             kQ_xR_x^4=C\sum_xR_x^3,\qquad
 \sum_xR_x^3\ge ck^{1+\gamma}.
\end{equation}
The latter inequality uses the integral in \eqref{eq:envelope}.

The beam vanishes on radius-two enlargements of these balls.
On one define
\[
 S_x^2=\sum_v|v(x)|^2,\qquad
 U_x(Y)=u(x+R_xY)/S_x,\quad |Y|\le2.
\]
Lemma~\ref{lem:scaled} gives all fixed-order squared derivative
bounds after this normalization, and the non-tail value fraction
is \(1-o(1)\). The value \(U_x(0)\) is a standard real Gaussian.

Let \(z\) be the coefficient vector with pair
\((\operatorname{Re}v(x),\operatorname{Im}v(x))\) for each packet,
and \(T\) the matrix of their four real \(Y\)-derivatives at zero.
Then \(|z|=S_x\), and the gradient covariance conditional on
the value is
\begin{equation}\label{eq:condcov}
 S_x^{-2}T\left(I-\frac{zz^{\mathsf t}}{S_x^2}\right)T^{\mathsf t}.
\end{equation}
The formula follows by an orthogonal decomposition of the
standard Gaussian coefficient vector along \(z\).
For each non-tail, the unit coefficient direction
\[
 c_v=(-\operatorname{Im}v(x),\operatorname{Re}v(x))/|v(x)|
\]
is perpendicular to \(z\); these directions are mutually
orthonormal, and
\[
 T c_v=|v(x)|\operatorname{Im}\frac{\dd_Yv}{v}(0).
\]
Equations~\eqref{eq:imaginary} and \eqref{eq:squares} give a
fixed positive trace lower bound in \eqref{eq:condcov}.
There is a deterministic unit direction \(e_x\) such that
\(Z_x=\partial_{e_x}U_x(0)\), conditional on \(U_x(0)\),
has variance between fixed positive constants. Unconditional
variance and the regression coefficient are bounded above
by the derivative square sums.

Those same estimates through order five, integrated on
\(B(0,2)\subset\R^4\), give
\begin{equation}\label{eq:c2moment}
 \E\|U_x\|_{C^2(\overline{B(0,1)})}^2
 \le C\,\E\|U_x\|_{H^5(B(0,2))}^2\le C.
\end{equation}
The Sobolev inequality here is the \(k=5,m=2,p=2,n=4\)
case of \cite[Theorem~3.49, p.~76]{Hunter}.
In particular, it controls the Euclidean gradient and
Hessian operator norms.

Take \(a'=1/4\), \(0<c'<a'/4\). For each sufficiently large
fixed \(H\), the conditional Gaussian density gives
\[
 \Prob(1\le Z_x\le2,\ |U_x(0)|<c'/H)\ge c_1/H.
\]
Its conditional mean is bounded on the indicated value interval,
and its conditional variance stays above and below fixed
positive numbers. Markov's inequality applied to
\eqref{eq:c2moment} costs at most \(C/H^2\) for the
failure of \(\|U_x\|_{C^2}\le H\).
Fix \(H\) large enough so the intersection has probability
\(\pi_*=c_1/(2H)>0\).

Taylor's formula at \(\pm a'e_x/H\) now gives opposite
values of absolute size at least \(a'/(2H)\).
The gradient bound keeps those signs on the parallel transverse
disks of radius \(\sigma_*=a'/(4H^2)\).
Their joining cylinder lies inside the unit ball.
The directions, disks and radii are deterministic; only success
depends on the Gaussian sample.

Let \(F_k\) be the fraction of successful balls weighted by
\(R_x^3\). Then \(0\le F_k\le1\), \(\E F_k\ge\pi_*\).
Writing \(q=\Prob(F_k\ge\pi_*/2)\) gives
\(\pi_*\le(\pi_*/2)(1-q)+q\), hence
\[
 \Prob(F_k\ge\pi_*/2)\ge\pi_* /(2-\pi_*)>0.
\]
No independence between balls is needed.
Each disk has three-area proportional to
\(\sigma_*^3R_x^3\). Summing and using \eqref{eq:packing}
proves the lemma.
\end{proof}

After all finite orders are fixed, Lemma~\ref{lem:signs}
has one positive probability \(p_*\), while failure in
Lemma~\ref{lem:jet} tends to zero.
For every sufficiently large allowed \(k\), probability
subtraction therefore gives a real sample satisfying
\eqref{eq:samplebounds}, the global first-jet lower bound,
and a certificate of size \(c_*k^{1+\gamma}\).

\subsection{Canceling the residual exactly}

\begin{lemma}\label{lem:repair}
Suppose the first-jet bound in Lemma~\ref{lem:jet} holds.
For an integer \(R\ge0\), assume the derivative bounds
\eqref{eq:samplebounds} for \(u\) through \(R+2\) and for
\(\mathfrak r\) through \(R+1\). There are smooth real corrections
to the pair, zero on every open set where \(\mathfrak r=0\), such that
\[
 \Div_\mu((A+\delta A)\dd u)+k^2(\rho+\delta\rho)u=0,\qquad
 \|(\delta A,\delta\rho)\|_{C^R}
       \le C_Rk^{-M+2129+1088R}.
\]
Thus \(M\ge2131+1088R\) gives corrections tending to zero in \(C^R\).
\end{lemma}

\begin{proof}
Put \(I=g_0^{-1}\), and define
\begin{equation}\label{eq:repair}
 \mathcal D=u^2+I(\dd u,\dd u),\quad
 F=-\mathfrak r/\mathcal D,\quad
 \delta A=FuI,\quad
 \delta\rho=k^{-2}\bigl(Fu-\Div_\mu(FI\dd u)\bigr).
\end{equation}
The first-jet bound implies
\(\mathcal D\ge \tfrac12 k^{-1000}W^2>0\).
Consequently all four expressions are globally smooth. The product rule
gives
\[
 \Div_\mu(FuI\dd u)
   +\bigl(Fu-\Div_\mu(FI\dd u)\bigr)u
     =F\bigl(u^2+I(\dd u,\dd u)\bigr)=-\mathfrak r.
\]
This is the required exact identity.

For \(0\le q\le R+1\), the bounds on \(u\) give
\[
 |D^q\mathcal D|\le C_q k^{2q+86}W^2.
\]
A derivative of order \(q\ge1\) of \(\mathcal D^{-1}\) is a sum
of terms
\(\mathcal D^{-m-1}\prod_{\nu=1}^m D^{q_\nu}\mathcal D\),
where \(q_\nu\ge1\), \(\sum q_\nu=q\), and \(m\le q\).
The lower bound for \(\mathcal D\) therefore yields, also at order zero,
\[
 |D^q\mathcal D^{-1}|
    \le C_q k^{1000+1088q}W^{-2},\qquad
 |D^qF|\le C_q k^{-M+1000+1088q}W^{-1}.
\]
These are estimates on already differentiated factors; they do not
differentiate \(W^{-1}\).
Multiplication by derivatives of \(u\) cancels the weights and gives
\[
 |D^q(Fu)|\le C_qk^{-M+1041+1088q},\qquad
 |D^q(FI\dd u)|\le C_qk^{-M+1043+1088q}.
\]
The divergence in \eqref{eq:repair} costs one derivative, followed
by \(k^{-2}\). Its exponent is
\(1043+1088(R+1)-2=2129+1088R\), as stated.
If \(\mathfrak r\) vanishes on an open set, then \(F\) and its
derivatives vanish there, so both corrections vanish there.
For sufficiently large \(k\) they preserve positivity, since the
input pair is positive on a compact manifold.
\end{proof}

\subsection{Making the eigenvalue simple while retaining its eigenfunction}

\begin{lemma}\label{lem:simple}
Suppose \(u\) is an exact real eigenfunction at \(k^2=N(N+3)\),
with \(N>1\). Suppose the pair is round and \(u=u_*\) on an open
region containing a point where \(0<t<1\) and \(\dd u,\dd t\)
are independent. A smooth tensor perturbation supported there can
make \(k^2\) simple while retaining \(u\) and \(\rho\).
Its size can be arbitrarily small in any prescribed finite \(C^R\) norm.
\end{lemma}

\begin{proof}
Take a small product coordinate patch \(O\) with coordinates
\((u,t,z_3,z_4)\). Choose a nonnegative bump \(\chi\) supported
inside \(O\), positive on a smaller product patch \(O_0\). On \(O\) put
\[
 E=\chi\left(I-
   \frac{(I\dd u)\otimes(I\dd u)}{I(\dd u,\dd u)}\right),
 \qquad I=g_0^{-1},
\]
and extend it by zero. The tensor is smooth, nonnegative and
\(E\dd u=0\). Let \(\mathcal E\) be the old real eigenspace at
\(k^2\). We first show that the kernel on \(\mathcal E\) of
\(\int_B E(\dd f,\dd f)\dd\mu\) is exactly \(\R u\).

If the integral is zero, positivity on \(O_0\) gives
\(\dd f\in\R\,\dd u\) there. On its product coordinates this
means \(f=b(u)\), with \(b\) smooth. The round equations imply
\begin{equation}\label{eq:composition}
 b''(u)|\dd u|_{g_0}^2=k^2\bigl(ub'(u)-b(u)\bigr).
\end{equation}
The ambient gradient of \(\operatorname{Re}(x_1+i x_2)^N\)
has squared norm \(N^2(1-t)^{N-1}\) and radial derivative \(Nu\).
Projection to the tangent space of the sphere gives
\[
 |\dd u|_{g_0}^2=N^2\bigl((1-t)^{N-1}-u^2\bigr).
\]
Differentiate \eqref{eq:composition} in the coordinate \(t\),
holding \(u\) fixed. Since \(N>1\) and \(t<1\), this gives
\(b''=0\); substitution eliminates the constant term of \(b\).
Thus \(f=cu\) on \(O_0\).
The difference satisfies a smooth uniformly elliptic second-order
equation. Aronszajn's theorem \cite[pp.~235--236]{Aronszajn}
applies: the smooth principal coefficients satisfy its
\(C^2\) and Lipschitz-second-derivative requirements, and the
lower-order coefficients are bounded locally. Open-set vanishing
gives infinite-order vanishing at an interior point. Unique
continuation, applied in overlapping charts on connected \(B\),
therefore gives \(f=cu\) globally.

Now replace \(A\) by \(A+\tau E\), with \(\tau>0\).
The same \(u\) still solves the exact equation at \(k^2\).
If that eigenvalue were not simple for arbitrarily small positive
\(\tau\), there would be \(\tau_m\downarrow0\) and corresponding
real eigenfunctions \(f_m\), of unit weighted \(L^2\) norm and
orthogonal to \(u\). Their energy identity
\[
 \int_B(A+\tau_mE)(\dd f_m,\dd f_m)\dd\mu=k^2
\]
bounds them in \(H^1\). Weak \(H^1\) compactness and compact
\(H^1\hookrightarrow L^2\) give a subsequence converging to
a norm-one \(f_0\) orthogonal to \(u\).
The weak equation passes to the limit; local elliptic regularity
makes \(f_0\) smooth, so \(f_0\in\mathcal E\).
For any fixed \(\phi\in\mathcal E\), testing the old equation
against \(f_m\) and the perturbed one against \(\phi\), then
subtracting and dividing by \(\tau_m\), gives
\[
 \int_B E(\dd f_m,\dd\phi)\dd\mu=0.
\]
Weak convergence preserves this identity. Taking \(\phi=f_0\)
places \(f_0\) in the kernel already shown to be \(\R u\),
which contradicts its norm and orthogonality.
Thus every sufficiently small positive \(\tau\) makes \(k^2\)
simple. The tensor \(E\) is fixed before this scalar is chosen,
so \(\tau E\) can meet any finite derivative tolerance.
The compactness and regularity used here are justified explicitly
in Lemma~\ref{lem:persistence} below.
\end{proof}

\subsection{Completion of the quantitative step}

\begin{proof}[Proof of Proposition~\ref{prop:step}]
Fix the pair, collar, \(r,\eta\). Choose \(R\ge\max(r,1)\) and
a smaller positive tolerance \(\eta'\) such that movement
less than \(\eta'\) in \(C^R\) implies the desired \(C^r\)
bound and keeps the pair positive and inside \(\mathcal U\).
This is possible because the fixed input is in the open neighborhood.
Set \(J=R+2\), take \(M\ge2131+1088R\), and then choose the
finite packet orders for \eqref{eq:samplebounds}.
These choices precede the frequency.

For every sufficiently large allowed \(k\), the simultaneous
sample above has a certificate of size \(c_*k^{1+\gamma}\)
and the hypotheses of Lemma~\ref{lem:repair}.
Choose \(k\) sufficiently large that its correction is smaller
than \(\eta'/2\) in \(C^R\).
The protected round region remains nonempty for every fixed
collar width: it suffices to require
\begin{equation}\label{eq:collarthreshold}
 Ck^{-1/3-100\delta}<w_{\rm c}/100,\qquad
 c w_{\rm c}k>(M'+P)\log k+\log C_q(w_{\rm c}).
\end{equation}
The second inequality makes a cutoff error
\(C_q(w_{\rm c})k^P e^{-c w_{\rm c}k}\) smaller than
\(k^{-M'}\). There are only finitely many derivative orders
and accuracy targets at this stage, so all these inequalities
hold for every sufficiently large \(k\).
Below \(s_1-w_{\rm c}/100\), \(u=u_*\) and the residual and
corrections vanish.

In \(l_-+w_{\rm c}/20<t<l_-+w_{\rm c}/10\), choose a point
where \(\sin(N\vartheta)\ne0\). Here \(\dd t\ne0\) and
the angular derivative of \(u_*\) is nonzero, so \(\dd u,\dd t\)
are independent. Apply Lemma~\ref{lem:simple}. After \(k\)
and the sample have been fixed, choose its scalar \(\tau\)
below the simplicity threshold with
\(\|\tau E\|_{C^R}<\eta'/2\).
All derivatives of \(E\) are finite at this stage; their possible
growth with \(k\) causes no obstruction to this later choice.
The perturbation is supported in \(t>l_-\).
It preserves the function, eigenvalue and signs exactly.

The final pair is round for \(t<l_-\), meets the norm tolerance,
and has a simple eigenvalue and the required certificate.
The number \(c_*>0\) is fixed before the universal choice of
large frequency. The exponent \(\gamma\), chart \(\Omega\)
and neighborhood \(\mathcal U\) were fixed earlier, independent
of all the finite orders, tolerances and collar widths.
\end{proof}
